\begin{figure}[htbp]
\centering
\begin{tikzpicture}[>=stealth,scale=1.3]
\draw[->,black!55] (-.2,0)--(3.1,0) node[right] {$e_1$};
\draw[->,black!55] (0,-.2)--(0,2.4) node[above] {$e_2$};
\draw[->,arkivblaa] (0,0)--(1,0) node[below] {$f_1$};
\draw[->,arkivblaa] (0,0)--(1,1) node[above left] {$f_2$};
\draw[->,black!80] (0,0)--(2,1) node[above] {$v$};
\draw[black!45,dashed] (0,1)--(2,1)--(2,0);
\node[align=left,anchor=west] at (3.5,1.2)
 {$S=\begin{pmatrix}1&1\\0&1\end{pmatrix}$\\[6pt]
 $[v]_e=\begin{pmatrix}2\\1\end{pmatrix}
 =S\begin{pmatrix}1\\1\end{pmatrix}=S[v]_f$};
\end{tikzpicture}
\caption{Samme vektor, to koordinatlister. Søjlerne i $S$ er de nye basisvektorer udtrykt i den gamle basis. Figuren bruger et ortonormalt $e$-system med samme dimensionsløse målestok; $f$ er ikke ortonormal. Et basisskift flytter ikke $v$.}
\label{fig:b1-basisskift}
\end{figure}
